\section{What the Composer Insists On}
\label{sec:ch11-what-the-composer-insists-on}
\index{composition!external@\code{@external} and requires@\code{@requires}|(}%

Chapter~\ref{ch:composition} produced its errors by taking the real pair of
schemas and applying one literal edit, on the grounds that an error message
from a toy schema is a fact about the toy. The three below are made the same
way, in \code{scripts/entity-cases.mjs}, and the messages are wgc 0.129.7's own
with the box it draws around them stripped off.

\subsection*{An \code{@external} field nothing requires}

Take the \code{@requires} off \code{shippingCost} and leave the
\code{@external} where it is:

\begin{minted}[fontsize=\footnotesize,breaklines]{text}
The subgraph "mosaic" could not be federated for the following reason:
The Object field "Product.category" is invalidly declared "@external". An Object field should only be declared
"@external" if it is part of a "@key", "@provides", or "@requires" field set, or the field is necessary to satisfy an
Interface implementation. In the case that none of these conditions is true, the "@external" directive should be
removed.
\end{minted}

The line breaks are the composer's own; it wraps to a fixed width inside its
table. Two lines above this one, wgc says \enquote{We found composition errors,
while composing} and invites you to check them below, which is where every
listing in this section starts and none of them repeats.

\index{composition!errors that name the fix}%
This is the most useful composition error in the book so far, and the reason is
that it enumerates the four legal cases and then tells you what to do. Compare
it with the three from chapter~\ref{ch:composition} that all reported as a
shareability problem on the key field and none of which named a key.

\index{suppress-warnings@\code{--suppress-warnings}}%
It also settles two flags chapter~\ref{ch:composition} could not exercise. That
chapter had nothing to say about \code{--suppress-warnings} because Mosaic
produced no warnings for it to suppress, and this was the first edit I made
that looked like a warning candidate. It is not one. The case runs the same
edit twice, once plain and once with the flag, and asserts that the two exit
codes match. They do, at 1. An orphaned \code{@external} is an error, and a
flag that suppresses warnings has nothing to say about it.

\index{ignore-external-keys@\code{--ignore-external-keys}}%
Chapter~\ref{ch:composition} handed the other one here by name:
\code{--ignore-external-keys} needs an \code{@external}, and Mosaic now has
one. Composing the committed pair with the flag produces a file byte-identical
to the one without it. That is the answer rather than an evasion, and the
reason is in the flag's name: Mosaic's \code{@external} field is not part of a
key. \code{Product} is keyed on \code{id}, which Mosaic resolves for itself.
Somebody whose key fields are marked \code{@external}, which is how federation
v1 required them to be written, would have something for the flag to ignore.

\subsection*{A \code{@requires} on a field you own}

The mirror image. Leave the \code{@requires} and take the \code{@external} off
\code{category}:

\begin{minted}[fontsize=\footnotesize,breaklines]{text}
The subgraph "mosaic" could not be federated for the following reason:
The field "Product.shippingCost" in subgraph "mosaic" defines a "@requires" directive with the following field set:
 "category".
However, neither the field "Product.category" nor any of its field set ancestors are declared "@external".
Consequently, "Product.category" is already provided by subgraph "mosaic" and should not form part of a "@requires"
directive field set.
\end{minted}

\index{requires@\code{@requires}!rules the composer enforces}%
Between them the two errors pin the pair: the directive and the declaration
only mean anything together. Cosmo's documentation states the three rules and
the composer enforces all three, which is worth saying because
chapter~\ref{ch:hotchocolate-quickly} found a documentation page whose sample
code does not compile and this book has been sceptical of vendor pages ever
since. The entity needs a \code{@key}; the fields named in \code{@requires}
must be \code{@external} here; and some other subgraph must actually
resolve them~\autocite{cosmo2026requires}.

\subsection*{The one that composed}

\index{composition!nullability of an @external copy}%
This is the case that changed my code, and it is the only one of the three that
produces no error at all.

Catalog declares \code{category: ProductCategory!}. My first version of the
\code{@external} copy in Mosaic was nullable, because the C\# property is
nullable and I let the schema follow the property without thinking about it.
The pair composes. Nothing is printed. And the client-facing schema comes out
like this:

\begin{minted}[fontsize=\footnotesize]{text}
catalog owns it and declares it     ProductCategory!
mosaic declares its @external copy  ProductCategory
the client-facing schema says       ProductCategory
the composer said                   nothing
\end{minted}

A field that Catalog guarantees became a field that might be null, for every
client of the whole graph, because a second service declared a copy it never
resolves and got the nullability wrong. There is no error, no warning, and
nothing in Mosaic's own schema that looks incorrect: \code{ProductCategory} is
a perfectly ordinary nullable enum field.

\index{external@\code{@external}!is a copy of somebody else's contract}%
The rule I take from it is that an \code{@external} declaration is not a
declaration at all. It is a copy of somebody else's, and it has to match theirs
in every respect that composition merges, which includes nullability and takes
the more permissive side. If your subgraph disagrees with the owner about a
type, composition believes whichever of you was more relaxed.

Hence the pair of attributes in
section~\ref{sec:ch11-a-field-from-the-other-service}:

\begin{minted}[fontsize=\footnotesize]{csharp}
[External]
[GraphQLNonNullType]
public ProductCategory? Category { get; set; }
\end{minted}

\code{[GraphQLNonNullType]} forces the schema to say what Catalog says. The
property stays nullable because it holds a value that arrives for some
representations and not others, and a non-nullable \code{ProductCategory} would
answer \code{FURNITURE} for absent, zero being a value of every C\# enum. Two
different nullabilities for the same thing, each right about its own layer.

\index{gates!a case for a behaviour that produces no error}%
All three are cases in the gate, and the third one is the argument for having a
gate at all. A composition error announces itself. A composed schema that is
subtly weaker than the one you meant announces nothing, and the only way I
noticed was reading the client schema the composer wrote. The case now reads it
for me on every run.

\medskip
So much for the directive that fetches something. The other half of the pair
promises not to.

\index{composition!external@\code{@external} and requires@\code{@requires}|)}%
